\documentclass[10pt,a4paper,onecolumn]{article}

\usepackage{amsmath,amssymb}
\usepackage{graphicx}
\usepackage[margin=1.5cm]{geometry}
\usepackage{hyperref}
\usepackage{cite}
\usepackage{enumitem}
\usepackage{booktabs}
\usepackage[bahasa]{babel}
\usepackage{titlesec}
\titlespacing*{\section}{0pt}{6pt}{3pt}
\titlespacing*{\subsection}{0pt}{4pt}{2pt}

\hypersetup{colorlinks=true,linkcolor=black,urlcolor=blue,citecolor=blue}
\setlength{\parskip}{0ex}
\setlength{\parindent}{0pt}
\setlength{\abovedisplayskip}{2pt}
\setlength{\belowdisplayskip}{2pt}
\setlength{\abovedisplayshortskip}{2pt}
\setlength{\belowdisplayshortskip}{2pt}

\title{\textbf{Deteksi Dini Deforestasi di Hutan Lindung \\
Menggunakan U-Net dan Citra Sentinel-2} \\
\small Idea Concept Paper --- Deforest.id}
\author{Tim Riset --- Deforest.id}
\date{}

\begin{document}
\maketitle
\thispagestyle{empty}

\paragraph{Abstrak.}
Kami mengusulkan sistem peringatan deforestasi otomatis untuk hutan
lindung Indonesia menggunakan model deep learning U-Net pada citra
Sentinel-2 resolusi 10~m. Inovasi utamanya adalah pipeline yang
sepenuhnya otonom: label pelatihan dihasilkan melalui NDVI temporal
differencing, menghilangkan kebutuhan anotasi manual. Kerangka kerja
peringatan dini tiga tingkat (Waspada/Siaga/Kritis) dengan amplifikasi
peringatan untuk kawasan lindung langsung terpetakan ke standar
pelaporan nasional dan internasional. Hasil awal menunjukkan IoU~0,73
dan Dice~0,84 pada data uji terpisah, dengan deteksi andal pada
0,64~ha per tile. Sistem ini dirancang untuk penyebaran cepat di
$\sim$94~juta hektar kawasan hutan Indonesia tanpa survei lapangan,
mendukung target FOLU Net Sink 2030.

\section{Latar Belakang}
Indonesia memiliki hutan tropis terbesar ketiga di dunia
94 juta ha, namun kehilangan 300 ribu ha hutan alam
setiap tahunnya \cite{fao2022sofo}. Hutan lindung sangat rentan
terhadap deforestasi skala kecil yang luput dari sistem monitoring
resolusi kasar. Petak $<$2~ha menyumbang $>$56\% emisi karbon hutan
tropis \cite{wahyuni2021dampak}, sementara fragmentasi mempercepat
hilangnya keanekaragaman hayati \cite{taubert2018fragmentation}.

Sistem operasional saat ini tidak menjembatani kesenjangan skala ini.
GFW beroperasi pada resolusi 30~m, melewatkan lahan di bawah 0,9~ha.
SIMONTANA Indonesia menggunakan minimum 6,25~ha untuk REDD+
\cite{unfccc2022indonesia} -- tiga kali lipat ambang 2~ha di mana
dampak ekologis menjadi parah. Tidak ada sistem yang menyediakan
\textit{prioritasi peringatan otomatis} untuk kawasan lindung sesuai
UU~No.~41/1999. Deforestasi pada skala 0,5--2~ha terlalu kecil untuk
sensor kasar dan terlalu besar untuk diabaikan secara ekologis
\cite{hansen2013high,nicfi2020satellite}.

\section{Metode}
Penelitian ini menggunakan 7 scene Sentinel-2 L1C dari Riau dan
Kalimantan Tengah (2019--2024), masing-masing mencakup hutan lindung
dengan tutupan awan $<$30\%, diproses dalam 5 kanal (B2, B3, B4, B8, QA60)
pada resolusi 10~m. Akuisisi melalui Google Earth Engine menghasilkan
60.587 tile $64\times64$ (50\% overlap): 40.484 baseline dan 20.103
deforestasi, dibagi menjadi 11.500 latih, 2.908 validasi, dan 5.695 uji.
Evaluasi menggunakan IoU, Dice, precision, dan recall pada data
hold-out \cite{powers2011evaluation,milletari2016vnet}.
Penelitian ini merupakan \textit{computational/design science research}
yang membangun pipeline akuisisi, pseudo-labeling, pelatihan, inferensi,
dan klasifikasi peringatan \cite{ronneberger2015unet,gorelick2017gee,tucker1979red,congalton1991accuracy}.

\section{Deskripsi Riset}
Pipeline menggabungkan citra Sentinel-2 (10~m, revisit 5 hari)
dengan U-Net untuk segmentasi pixel-level deforestasi. Inovasi inti
adalah NDVI temporal differencing yang menghasilkan label pelatihan
tanpa anotasi manual \cite{tucker1979red}:
\begin{equation}
\Delta\text{NDVI}(x,y) = \text{NDVI}_{\text{sebelum}} - \text{NDVI}_{\text{sesudah}}
\quad\Longrightarrow\quad
\text{mask}(x,y) = 
\begin{cases}
1 & \Delta\text{NDVI} < -0,15 \\
0 & \text{lainnya}
\end{cases}
\end{equation}
ditutup morfologis ($3\times3$) untuk menghilang derau. Ambang $-0,15$
dikalibrasi dari 50 sampel acak. Pendekatan ini memungkinkan
skalabilitas lintas provinsi tanpa kampanye lapangan
\cite{gonzalez2008digital,goodfellow2016deep}.

Area deforestasi diklasifikasikan ke tiga tingkat: \textit{Waspada}
($\ge 0,64$~ha, setara minimum hutan FAO), \textit{Siaga}
($\ge 2$~ha, ambang fragmentasi ekologis), dan \textit{Kritis}
($\ge 6,25$~ha, minimum REDD+ Indonesia). Untuk hutan lindung,
peringatan dinaikkan satu tingkat \cite{esa2024sentinel2,renno2016amazon,ri1999kehutanan}.

Deteksi awan menggunakan heuristik
$R{>}180 \land G{>}180 \land B{>}180 \land \text{NDVI}{<}0,10$;
sisa noise ditangani Dice loss \cite{zhu2012cloud,milletari2016vnet}.

Pipeline beroperasi dalam enam tahap
\cite{kingma2015adam,shapiro2001computer}:
akuisisi 7 scene via GEE, tiling $64\times64$ overlap~50\%
(60.587 tile), cloud masking, pseudo-labeling, pelatihan U-Net
(4 encoder/decoder, 32--512 filter, input 5 kanal, Dice loss, Adam
$1\times10^{-4}$, batch 64, 50 epoch), dan inferensi sliding window
dengan connected-component labeling serta agregasi area per ambang.

\section{Hasil Awal, Dampak, dan Skalabilitas}
Pelatihan pada RTX~3060 (12~GB) selesai ${<}$3~jam (50 epoch).
Pada 5.695 tile uji dari 2 scene tak terlihat \cite{powers2011evaluation}:
IoU~0,7256, Dice~0,8410, presisi~0,8340, recall~0,8480. Sistem
andal mendeteksi deforestasi pada 0,64~ha; sebagai perbandingan,
NDVI thresholding sederhana hanya mencapai $\sim$0,53~F1,
mengonfirmasi U-Net mempelajari konteks spasial
\cite{goodfellow2016deep}.

Sistem memberikan peringatan harian otomatis pada resolusi 0,01~ha --
peningkatan 900$\times$ dari minimum REDD+ 6,25~ha. Biaya anotasi nol
(pseudo-labeling dari pasangan Sentinel-2, $\sim$30 menit cloud per
provinsi). Pelatihan ${<}$3~jam pada GPU konsumen (\$250), inferensi
${<}$5~menit per scene -- tanpa infrastruktur cloud khusus. Ambang
peringatan terpetakan langsung ke REDD+, FAO, dan UU~No.~41/1999,
mendukung target FOLU Net Sink 2030
\cite{moef2022folu,gorelick2017gee,abadi2016tensorflow,ipcc2019land}.

\section{Kesimpulan}
Deteksi dini deforestasi di hutan lindung Indonesia dapat dilakukan
secara kuantitatif melalui kombinasi Sentinel-2 dan U-Net dengan
pipeline otonom (pseudo-labeling NDVI, segmentasi, cloud masking,
klasifikasi peringatan). Hasil awal menjanjikan (IoU~0,7256,
Dice~0,8410). Dengan ambang 0,64/2/6,25~ha, sistem berpotensi menjadi
dasar pemantauan deforestasi dini yang lebih cepat, murah, dan selaras
dengan target FOLU Net Sink 2030.

\small
\renewcommand{\refname}{Daftar Pustaka}
\bibliographystyle{ieeetr}
\bibliography{icp-refs}
\end{document}
